\section*{Course Structure and Chapter Roadmap}
\addcontentsline{toc}{section}{Course Structure and Chapter Roadmap}

The course is organized into twelve principal chapters, corresponding approximately to twelve teaching weeks. The chapters are not independent units. Each one introduces tools that are used later, and the order is chosen so that the formalism grows from familiar mechanics toward relativistic fields, electrodynamics, propagation, radiation, and geometric field theories.

The first three chapters establish the field concept and the variational method. Chapters 4--6 introduce the relativistic language, the scalar-field prototype, and the relation between symmetry and conservation. Chapters 7--9 develop electromagnetism as a complete classical field theory, including gauge freedom and energy-momentum. Chapters 10--11 explain how fields generated by sources are solved and how radiation emerges. Chapter 12 gives a geometric outlook toward gauge theory and gravity.

\subsection*{Chapter 1: Fields, Locality, and Continuous Degrees of Freedom}

The opening chapter establishes what a field is and why fields are required. We begin with familiar examples: temperature as a scalar field, the velocity of a fluid as a vector field, the displacement of a string as a time-dependent field, and the electric field as a physical quantity defined throughout space. These examples make clear that a field assigns data to every point of a domain and therefore carries infinitely many degrees of freedom.

The idea of locality is then introduced. A local field theory describes changes at a point in terms of the field and its derivatives near that point. This contrasts with a direct action-at-a-distance description. The distinction will be illustrated through sources, responses, and propagation. We will ask what it means for a source to affect a field locally and how information can move through a continuous medium.

A central worked construction will be the passage from a discrete chain of coupled oscillators to a continuous string. The discrete labels will be replaced by a spatial coordinate, finite differences will become derivatives, and the sum over masses will become an integral. This example will show explicitly how a field theory can emerge as the continuum limit of a many-particle system.

By the end of the chapter, the student should be able to identify scalar, vector, and tensor fields, state their domains and values, distinguish discrete from continuous degrees of freedom, and explain why local differential equations are natural equations of motion for fields.

\subsection*{Chapter 2: Action Principles: From Particles to Fields}

The second chapter develops the variational principle in the setting of ordinary mechanics. This is necessary because the field-theoretic action is a direct extension of the mechanical action. We introduce a functional as an object whose input is an entire trajectory rather than a single number. The action is then defined as an integral of the Lagrangian along a path.

The variation of a trajectory will be constructed explicitly. We will compare a physical trajectory with nearby trajectories that share fixed endpoints, differentiate the action with respect to the variation parameter, and use integration by parts. The boundary term will not simply be discarded; the endpoint conditions that make it vanish will be stated and interpreted.

The Euler--Lagrange equations will be derived in full and applied to several concrete systems: a free particle, a particle in a potential, the harmonic oscillator, and a simple constrained coordinate. These examples will show how to move from a Lagrangian to an equation of motion and then verify that a proposed solution satisfies the resulting equation.

This chapter also establishes a working habit that will recur throughout the course: whenever an equation is derived from an action, the allowed variations, the boundary data, and the variables held fixed must be specified.

\subsection*{Chapter 3: Lagrangian Field Theory and Field Equations}

The third chapter performs the transition from trajectories to fields. The Lagrangian becomes a Lagrangian density, and the action becomes an integral over spacetime. A variation now changes the field value at every point, and derivatives of the field produce derivatives of the variation.

The Euler--Lagrange field equations will be derived carefully for one field and then generalized to several coupled fields. The calculation will display the spacetime integration by parts and the boundary term on the boundary of the integration region. This provides the basic method used in nearly every later derivation.

Several field equations will be obtained from actions. The vibrating string will produce the wave equation. A static energy functional will lead to Laplace or Poisson equations. A field with a potential \(V(\phi)\) will provide a first example of a nonlinear field equation. We will compare the roles of kinetic, gradient, source, and potential terms in a Lagrangian density.

The chapter will close by identifying canonical field momenta and the Hamiltonian density at an introductory level. The main goal, however, is practical: given a Lagrangian density, the student should be able to vary it, manage the boundary term, and derive the corresponding differential equation.

\subsection*{Chapter 4: Spacetime, Lorentz Symmetry, and Covariant Notation}

Classical relativistic field theory requires a language in which space and time are combined consistently. This chapter introduces Minkowski spacetime, events, the spacetime interval, proper time, and Lorentz transformations. The metric convention will be fixed and used throughout the remaining chapters.

Four-vectors, covectors, and tensors will be introduced through calculations rather than definitions alone. Students will practice raising and lowering indices, contracting tensors, distinguishing free and dummy indices, and checking whether an expression is a Lorentz scalar. Four-position, four-velocity, four-momentum, and the four-gradient will provide the main examples.

The d'Alembert operator will be constructed from the metric and the four-gradient. Familiar equations such as the wave equation will then be rewritten in covariant form. We will compare the compact tensor notation with the expanded time-and-space components so that the notation does not conceal the underlying differential equation.

Particular attention will be given to sign conventions. The notes will show how the chosen metric signature affects norms, derivatives, and Lagrangian densities. By the end of the chapter, the student should be able to read and manipulate the index notation used in the remainder of the course.

\subsection*{Chapter 5: The Real Scalar Field}

The real scalar field is the simplest complete relativistic field theory and will serve as the principal training model. We construct the free scalar-field Lagrangian from Lorentz scalars and derive the massless and massive field equations. The massive equation is the classical Klein--Gordon equation, treated here as a classical partial differential equation rather than as a quantum equation.

Plane-wave solutions will be substituted explicitly into the field equation. This calculation will produce the dispersion relation and clarify the meaning of frequency, wave vector, and the mass parameter. The massless and massive cases will be compared, including phase and group behavior where appropriate.

The chapter will then examine energy and momentum for scalar fields. We will calculate the energy density of simple configurations and ask under what conditions the energy is bounded below. Interacting theories with a potential \(V(\phi)\) will be introduced, and equilibrium configurations will be identified as extrema of the potential.

Linearization around an equilibrium will show how a nonlinear theory produces an approximate linear wave equation for small perturbations. This provides a useful method that reappears in many areas of physics: find a background solution, perturb it, keep the leading terms, and analyze stability.

\subsection*{Chapter 6: Symmetries, Currents, and Noether's Theorem}

This chapter develops one of the central organizing principles of field theory: continuous symmetries produce local conservation laws. We begin with infinitesimal transformations of coordinates and fields, making clear what changes actively and what changes because the coordinates have been relabeled.

Noether's theorem will be derived for fields with enough detail to identify every term in the current. The derivation will distinguish variations that vanish at the boundary from transformations that represent physical symmetries. The resulting divergence equation will be related to a conserved charge by integration over space.

Several explicit applications will follow. Time translation will be connected with energy, spatial translation with momentum, and rotations with angular momentum. A complex scalar field will provide an example of an internal global phase symmetry and its associated current. For each case, the current will be calculated and its conservation checked using the field equation.

The chapter will also prepare the transition to gauge theory by distinguishing a global symmetry from a local redundancy. The student should leave with a practical algorithm: specify the infinitesimal transformation, compute the variation of the Lagrangian density, identify the Noether current, and verify its divergence.

\subsection*{Chapter 7: Electromagnetic Potentials and Gauge Freedom}

Electromagnetism is introduced through potentials because the potentials are the natural variables of the action principle. We begin with the scalar and vector potentials and derive the electric and magnetic fields from them. The relations will be checked against the homogeneous Maxwell equations.

Gauge transformations will then be performed explicitly. The potentials change, but the electric and magnetic fields remain unchanged. This calculation provides the first concrete example of a description containing more variables than there are physical degrees of freedom. The distinction between a change of description and a change of physical state will be emphasized.

The scalar and vector potentials will be combined into the four-potential \(A_{\mu}\). From it we construct the antisymmetric electromagnetic field tensor \(F_{\mu\nu}\). The matrix of components will be written out, and the electric and magnetic components will be identified with careful attention to signs and index positions.

The two standard Lorentz invariants of the electromagnetic field will be calculated and interpreted. Finally, the homogeneous Maxwell equations will be written in covariant form and related to the fact that the field tensor is constructed from a potential.

\subsection*{Chapter 8: Maxwell Theory from an Action}

This chapter builds Maxwell theory as a variational field theory. The electromagnetic Lagrangian density is formed from the field tensor, and a coupling term to an external four-current is added. We then vary the action with respect to the four-potential.

The calculation will be shown in full: the variation of the field tensor, the use of antisymmetry, the integration by parts, the treatment of the boundary term, and the extraction of the inhomogeneous Maxwell equations. The covariant result will be expanded into Gauss's law and the Ampere--Maxwell law.

Gauge invariance of the source coupling will be analyzed. We will show that the action changes by a boundary term only when the current is conserved. This makes the relation between gauge invariance and charge conservation explicit rather than merely asserted.

The Lorenz gauge condition will be introduced and used to reduce the equations for the potentials to wave equations. The action for a relativistic charged particle interacting with the electromagnetic potential will then be varied to obtain the Lorentz-force law. This completes a unified action-based description of both the field and its interaction with matter.

\subsection*{Chapter 9: Energy, Momentum, and Stress in Field Theory}

Field equations describe evolution, but they do not by themselves show how energy and momentum are distributed and transported. This chapter develops the local energy-momentum description of scalar and electromagnetic fields.

We begin with canonical field momentum and Hamiltonian density. Translation symmetry is then used to construct the canonical energy-momentum tensor. The distinction between canonical and symmetric tensors will be discussed, including why the symmetric form is especially natural in relativistic theories and in coupling to gravity.

For the scalar field, the energy density, momentum density, and flux will be computed. For the electromagnetic field, the components of the energy-momentum tensor will be identified with electromagnetic energy density, the Poynting vector, momentum density, and the Maxwell stress tensor.

Several explicit calculations will show how these quantities are used. We will calculate energy flow in a plane wave, force from the stress tensor in a simple geometry, and radiation pressure. Local conservation equations will be integrated over a finite region to obtain balance laws with flux through the boundary.

\subsection*{Chapter 10: Green Functions, Sources, and Field Propagation}

The tenth chapter turns from deriving field equations to solving them. Green functions provide a systematic method for linear differential equations with sources. We introduce the idea through an operator equation and define a Green function as the response to a point source.

The static case will be developed first. The Green function of the Poisson equation in three-dimensional space will be used to recover the Coulomb potential. The solution for a distributed source will then be obtained by superposition. This calculation will connect the abstract Green-function definition with a familiar physical result.

Time-dependent equations require additional choices. We will construct retarded and advanced Green functions for wave equations and explain the boundary conditions encoded by each choice. The retarded solution will be associated with causal propagation from past sources to later fields.

The method will then be applied to electromagnetic potentials. Retarded potentials will be derived as integrals over the source evaluated at retarded time. The chapter will emphasize that a differential equation alone does not determine a unique physical solution; boundary and causal conditions are part of the problem.

\subsection*{Chapter 11: Radiation from Time-Dependent Sources}

Radiation is the regime in which a time-dependent source creates a field that carries energy to arbitrarily large distances. We begin by separating near-field terms from radiation-field terms according to their dependence on distance. This makes clear why a \(1/r\) field can transport a finite amount of energy through a large sphere.

A multipole expansion will organize the radiation produced by localized sources. The electric dipole will be studied in detail. Starting from an oscillating dipole moment, we will calculate the far-zone fields, the Poynting vector, the angular distribution of emitted power, and the total radiated power.

The Lienard--Wiechert potentials will then be introduced for a moving point charge. Their dependence on retarded time will be analyzed, and the velocity and acceleration contributions to the field will be distinguished. The acceleration field will be identified as the radiative part.

The Larmor formula will be derived or motivated from the radiation field, with its assumptions stated clearly. Energy balance and radiation reaction will be discussed at an introductory level, including the conceptual difficulties that arise when a point charge interacts with its own field.

\subsection*{Chapter 12: Gauge Fields and Gravity: The Geometric Outlook}

The final chapter places the preceding material in a broader geometric framework. Electromagnetism will be reinterpreted as a \(U(1)\) gauge theory. The four-potential will be viewed as a connection, the field tensor as curvature, and the gauge-covariant derivative as the derivative appropriate for charged fields.

This language will be introduced through the electromagnetic case before any non-Abelian generalization is attempted. The non-Abelian field strength will then be presented as a natural extension in which the gauge potentials themselves interact because the symmetry group is noncommutative.

Gravity will be introduced as a theory in which the metric is a dynamical field. Connections, curvature, geodesics, and the Einstein--Hilbert action will be discussed at a structural level. The variation of the gravitational action will be outlined sufficiently to show why the Einstein tensor appears and why stress-energy acts as a source.

The weak-field form \(g_{\mu\nu}=\eta_{\mu\nu}+h_{\mu\nu}\) will connect general relativity with linear field theory and the Newtonian potential. The chapter will conclude by comparing electromagnetism and gravity: their dynamical variables, symmetries, sources, field strengths, conservation laws, and geometric interpretations. The purpose is not to complete a course in Yang--Mills theory or general relativity, but to show where the methods developed in the previous eleven chapters lead.
